\documentclass[11pt]{article}
\input{packages.tex}
\DeclareMathOperator{\sign}{sign}
\DeclareMathOperator{\Imm}{Im}

% Настройка колонтитулов
\pagestyle{fancy}
\fancyhf{} % Очистить все поля
\fancyhead[L]{Ах, Дюдюка! Ух, Дюдюка! Бяка, привередина}
\fancyhead[R]{\href{https://t.me/sol_m_07}{Соля}}
\fancyfoot[C]{\thepage}               % Справа внизу

\setlength{\headheight}{13.6pt}

\hypersetup{
    colorlinks=true,
    linkcolor=blue,
    citecolor=blue,
    filecolor=blue,
    urlcolor=blue,
}

\begin{document}

\begin{titlepage}
    \centering
    \vspace*{\fill}  % заполняет пространство сверху

    {\Huge Лекция 2. Критерий Поста. Формальные системы
    \par}

    \vspace*{\fill}  % заполняет пространство снизу
\end{titlepage}

\setcounter{page}{1}

\section{Классы Поста (из прошлой лекции)}

\paragraph{Функции, сохраняющие ноль}
\[T_0 = \{f(x_1, \ldots, x_n) \mid f(0, \ldots, 0) = 0\}\]
\paragraph{Функции, сохраняющие единицу}
\[T_1 = \{f(x_1, \ldots, x_n) \mid f(1, \ldots, 1) = 1\}\]
\paragraph{Линейные функции}
\[L \text{ --- множество всех таких функций } f(x_1, \ldots, x_n) \text{, что для них}\]
\[\text{полином Жегалкина имеет не более, чем первую степень}\]
\paragraph{Самодвойственные функции}
\[S = \{f(x_1, \ldots, x_n) \mid f(\overline{x_1}, \ldots, \overline{x_n}) = \overline{f(x_1, \ldots, x_n)}\}\]
\paragraph{Монотонные функции}
\[M = \{f(x_1, \ldots, x_n) \mid \alpha \succeq \beta \implies f(\alpha) > f(\beta)\}\]
Здесь $"\succeq"$ --- частичный порядок, заданный на множестве всех $n$-ок из нулей и единиц. При этом 
\[\alpha \succeq \beta \iff \alpha_i \geq \beta_i \quad \forall i \in \{1, \ldots, n\}\]

\section{Критерий Поста}

\begin{theorem}[Критерий Поста]
    Множество функций полно тогда и только тогда, когда оно не содержится 
    целиком ни в одном из классов Поста.
\end{theorem}

\begin{proof}
    $(\Longrightarrow)$: Пусть множество функций $\mathcal{E}$ полно. \\ 
    Учтём тот факт, что все классы Поста замкнуты \textbf{($\ast$)} и не совпадают с множеством 
    всех функций алгебры логики \textbf{($\ast\ast$)} (автор докажет это, когда будет время). \\ 
    От противного предположим, что $\mathcal{E}$ содержится в одном из 
    классов Поста. \\ Из \textbf{($\ast$)} получаем, что из $\mathcal{E}$ 
    можно получить только функции из того же класса Поста. \\ 
    Тогда из \textbf{($\ast\ast$)} получаем, что замыкание $\mathcal{E}$ 
    не совпадает со множеством всех функций алгебры логики $\implies \mathcal{E}$ 
    не полно -- противоречие. \\ \\
    $(\Longleftarrow)$: Пусть множество функций $\mathcal{E}$ не содержится 
    целиком ни в одном из классов Поста. Распишем это условие в следующем виде: \\
    \[\exists f_{\overline{0}} \notin T_0, \ f_{\overline{1}} \notin T_1, \ 
    f_{\overline{S}} \notin S, \ f_{\overline{L}} \notin L, \ f_{\overline{M}} \notin M\] 

    \begin{remark}
        Эти функции (по умолчанию лежащие в $\mathcal{E}$) могут совпадать. 
        Это ничего не изменит. Главное, заметить, что переписанное условие равносильно исходному.
    \end{remark}

    \paragraph{Шаг 1}
    Рассмотрим $f_{\overline{0}}(x_1, x_2, ..., x_k)$. Отождествим все переменные с 
    иксом и получим $f_{\overline{0}}(x, x, ..., x)$, что можно представить, как 
    функцию от одного переменного $\varphi_{\overline{0}}(x)$. Вспоминаем, что 
    $f_{\overline{0}} \notin T_0 \implies f_{\overline{0}}(0, 0, ..., 0) = 
    \varphi_{\overline{0}}(0) = 1$. Тогда \[\varphi_{\overline{0}}(x) = 1 
    \text{ или } \varphi_{\overline{0}}(x) = \overline{x}\] \\ Аналогично, для 
    $f_{\overline{1}}(x_1, x_2, ..., x_k)$ получим ситуацию $f_{\overline{1}}(1, 1, 
    ..., 1) = \varphi_{\overline{1}}(1) = 0$. Тогда \[\varphi_{\overline{1}}(x) = 0 
    \text{ или } \varphi_{\overline{1}}(x) = \overline{x}\]
    \paragraph{Шаг 2}
    Рассмотрим все варианты продолжения событий:
    \subparagraph{Вариант 1} 
    Пусть получили $\overline{x}$ и константу. Тогда путём суперпозиции получим 
    вторую константу ($\overline{0} = 1 \text{ или } \overline{1} = 0$).
    \subparagraph{Вариант 2} 
    Пусть получили только $\overline{x}$. Рассмотрим $f_{\overline{S}}$. Так как 
    $f_{\overline{S}} \notin S$, то \[\exists \alpha \text{ - набор переменных: } 
    f_{\overline{S}}(\alpha) = f_{\overline{S}}(\overline{\alpha})\] 
    Определим "степень": \[x^\sigma = \begin{cases}
        x, \ \sigma = 1 \\
        \overline{x}, \ \sigma = 0
    \end{cases}\]
    Тогда можем образовать функцию от одного переменного: 
    \[f_{\overline{S}}(x^{\alpha_1}, x^{\alpha_2}, ..., x^{\alpha_k}) = 
    \varphi_{\overline{S}}(x)\]
    Получаем \[\varphi_{\overline{S}}(x) = \varphi_{\overline{S}}(\overline{x}) 
    \implies \varphi_{\overline{S}}(x) = 0 \text{ или } \varphi_{\overline{S}}(x) = 1\] 
    \subparagraph{Вариант 3}
    Пусть получили 0 и 1. Рассмотрим $f_{\overline{M}}$. Так как 
    $f_{\overline{M}} \notin M$, то \[\exists \alpha \succcurlyeq \beta: 
    f_{\overline{M}}(\alpha) = 0 \text{ и } f_{\overline{M}}(\beta) = 1 \ \textbf{($\ast$)}\]
    От набора $\beta$ к набору $\alpha$ можно перейти, последовательно меняя нули 
    на единицы (тогда соседние наборы будут различаться только в одном элементе). 
    Порядок изменения нам не важен. Пример:

    \begin{figure}[h]
    \centering
    \includegraphics[width=0.3\linewidth]{pictures/Переход от одного набора к другому.png}
    \end{figure} \leavevmode
    \\
    Заметим, что из \textbf{($\ast$)} следует, что $\exists \widetilde{\alpha}, 
    \widetilde{\beta}: (\widetilde{\alpha} \succcurlyeq \widetilde{\beta}) 
    \wedge (f_{\overline{M}}(\widetilde{\alpha}) = 0 \text{ и } 
    f_{\overline{M}}(\widetilde{\beta}) = 1)$, где $\widetilde{\alpha}$ и 
    $\widetilde{\beta}$ -- соседние наборы в переходе выше. Раз они различаются 
    только в одном элементе, то отождествим его с иксом и получим функцию от одного 
    переменного: \[f_{\overline{M}}(\widetilde{\alpha}_1, ..., \widetilde{\alpha}_{i - 1}, 
    x, \widetilde{\alpha}_{i + 1}, ..., \widetilde{\alpha}_k) = \varphi_{\overline{M}}(x)\]
    Получаем: \[\begin{cases}
        \varphi_{\overline{M}}(0) = 1 \\
        \varphi_{\overline{M}}(1) = 0
    \end{cases} \implies \ \varphi_{\overline{M}}(x) = \overline{x}\]

    \paragraph{Шаг 4}
    Во всех вариантах в итоге получили набор функций $\overline{x}, 0, 1$. Теперь 
    рассмотрим $f_{\overline{L}}$. Раз она не линейна, то в полиноме Жегалкина 
    найдётся слагаемое степени $> 1$. Тогда рассмотрим слагаемое наименьшей степени 
    в \textbf{нелинейной части} полинома Жегалкина (можно воспринимать как самое левое). 
    Оно будет вида $x_{i_1}x_{i_2}...x_{i_t}$. 
    \begin{enumerate}
        \item Отождествим $x_{i_1}$ с $x$;
        \item Отождествим $x_{i_2}$ с $y$;
        \item Заединичим $x_{i_3},...,x_{i_t}$ (ничего не сделаем, если иксов всего два);
        \item Остальные (не входящие в наше слагаемое) занулим.
    \end{enumerate}
    Получим $\varphi_{\overline{L}}(x, y) = 1 + x + y + xy$, где каждое из первых трёх 
    слагаемых может не входить. \\
    Избавимся от всего лишнего: 
    \begin{enumerate}
        \item Если есть 1, то рассмотрим $\lnot \varphi(x, y)$, так как 
        $\varphi(x, y) + 1 = \urcorner \varphi(x, y)$
        \item Если есть $x$, то рассмотрим $\varphi(x, \overline{y})$, так как 
        $x + xy = x(1 + y) = x \overline{y}$
        \item Аналогично если есть $y$, то рассмотрим $\varphi(\overline{x}, y)$
    \end{enumerate}
    Таким образом мы избавимся от всего и оставим только конъюнкцию. Тогда имеем 
    конъюнкцию и отрицание $\implies$ можем получить любую функцию алгебры логики 
    $\implies$ множество функций полно.
\end{proof}

\section{Формальные системы}

\begin{definition}
    Формальная система -- это набор из четырёх множеств $\{\Sigma, \mathcal{F}, 
    \mathcal{A}, \mathcal{R}\}$, где \begin{itemize}
        \item $\Sigma$ -- алфавит (счётное множество);
        \item $\mathcal{F}$ -- формулы ($\mathcal{F} \subseteq \Sigma^*$ -- все 
        конечные слова в алфавите $\Sigma$);
        \item $\mathcal{A}$ -- аксиомы ($\mathcal{A} \subseteq \mathcal{F}$);
        \item $\mathcal{R}$ -- правила вывода (о них чуть поподробнее)
    \end{itemize}
    Пусть $\mathcal{R}_i \in \mathcal{R}$. Для него зафиксировано $k_i \in \mathbb{N}:$ 
    \[\mathcal{R}_i \colon \mathcal{F}^{k_i} \to \mathcal{F}\] 
    При этом $\mathcal{R}_i$ -- частично заданная функция, т.е. некоторой формуле 
    длины $k_i$ из $\mathcal{F}^{k_i}$ может быть не сопоставлена формула.
    \begin{example}
        Можем преобразовать \ $(x + 1)(x + y) = 25$ \ в \ $x^2 + xy + x + y = 25$ \ 
        (здесь мы работаем с синтаксисом)
    \end{example}
\end{definition}

\subsection{Исчисление высказываний}

\paragraph{Алфавит ($\Sigma$)} 
\begin{enumerate}
    \item счётное количество символов переменных
    \item $\lnot, \ \to, \ (, \ )$
\end{enumerate}

\paragraph{Формулы ($\mathcal{F}$)}
\begin{itemize}
    \item Если $A$ -- символ переменной, то $A$ -- формула.
    \item Если $A = (\lnot B)$, где $B$ -- формула, то $A$ -- формула.
    \item Если $A = (B \to C)$, где $B$ и $C$ -- формулы, то $A$ -- формула.
    \item Других формул нет (не всегда уточняют, но мы это делаем).
\end{itemize}

\paragraph{Аксиомы ($\mathcal{A}$)} 
Аксиом бесконечно много, поэтому зададим их схемами:
\begin{enumerate}
    \item $(A \to (B \to A))$
    \item $((A \to (B \to C)) \to ((A \to B) \to (A \to C)))$
    \item $(((\lnot B) \to (\lnot A)) \to (((\lnot B) \to A) \to B))$
\end{enumerate}
Вместо $A, B$ и $C$ можно подставлять любые формулы, но для одинаковых букв 
должны быть одинаковые формулы.

\paragraph{Правила вывода ($\mathcal{R}$)} 
modus ponens $A, (A \to B) \vdash B$ \\ \\
Автор немного почитал в интернете про эту непонятную темку и переводит выше написанное: \\ \\
\textbf{modus ponens} (с латыни "правило отделения" \ или "утверждающий модус") -- правило, которое 
говорит нам о том, что истинность $A$ и истинность $(A \to B)$ влекут за собой истинность $B$ 
($\vdash$ -- знак "выводится"). \\ \\
На лекции было оговорено, что с помощью \textbf{modus ponens} можно получать тавтологии 
(формулы, истинные при любых значениях входящих в них переменных). Действительно, 
заметим, что если $A$ и $(A \to B)$ -- тавтологии, то и $B$ является тавтологией.
\\ \\ \\
\begin{figure}[h]
    \centering
    
\includegraphics[width=1\linewidth]{pictures/Дюдюка.png}
\end{figure}

\end{document}